\documentclass{article}
\usepackage{amsmath,amssymb,amsthm}
\usepackage{algorithm,algorithmic}
\usepackage{bm}
\usepackage{hyperref}
\newtheorem{theorem}{Theorem}
\newtheorem{lemma}{Lemma}
\newtheorem{assumption}{Assumption}
\newcommand{\prox}{\operatorname{prox}}
\newcommand{\grad}{\nabla}
\newcommand{\R}{\mathbb{R}}
\newcommand{\EE}{\mathbb{E}}

\begin{document}

\section*{Proximal Gradient Method (PGM)}

\section*{Mathematical Formulation}
Let 
\[
F(x)\;:=\;f(x)+g(x),\qquad x\in\R^{d},
\]
where  

\begin{itemize}
  \item $f:\R^{d}\rightarrow\R$ is convex and differentiable with $L$‑Lipschitz continuous gradient,
        \[
        \|\grad f(x)-\grad f(y)\|\le L\|x-y\|,\qquad\forall x,y\in\R^{d};
        \]
  \item $g:\R^{d}\rightarrow\R\cup\{+\infty\}$ is convex (possibly non‑smooth) and proper.
\end{itemize}
For a stepsize $\eta>0$ the proximal gradient update reads
\begin{equation}
\label{eq:update}
x_{k+1}\;=\;\prox_{\eta g}\!\bigl(x_{k}-\eta\,\grad f(x_{k})\bigr),
\qquad k=0,1,2,\dots
\end{equation}
where the proximal operator of $\eta g$ is defined by
\[
\prox_{\eta g}(v)\;:=\;\arg\min_{u\in\R^{d}}\Bigl\{g(u)+\frac{1}{2\eta}\|u-v\|^{2}\Bigr\}.
\]

\section*{Pseudocode}
\begin{algorithm}[H]
\caption{Proximal Gradient Method}
\begin{algorithmic}[1]
\REQUIRE Initial point $x_{0}\in\R^{d}$, stepsize $\eta\in(0,1/L]$.
\FOR{$k=0,1,\dots,T-1$}
    \STATE Compute gradient $g_{k}\gets \grad f(x_{k})$.
    \STATE Take a gradient step $y_{k}\gets x_{k}-\eta g_{k}$.
    \STATE Apply proximal map $x_{k+1}\gets \prox_{\eta g}(y_{k})$.
\ENDFOR
\RETURN $x_{T}$
\end{algorithmic}
\end{algorithm}

\section*{Dry Run Example}
Consider the scalar problem $f(w)=(w-3)^{2}$, $g\equiv0$, stepsize $\eta=0.1$, and $w_{0}=0$.

\begin{center}
\begin{tabular}{c|c|c|c}
$k$ & $\grad f(w_{k})$ & $w_{k+1}=w_{k}-\eta\grad f(w_{k})$ & $F(w_{k+1})$\\\hline
0 & $2(w_{0}-3)=-6$ & $0-0.1(-6)=0.6$ & $(0.6-3)^{2}=5.76$\\
1 & $2(0.6-3)=-4.8$ & $0.6-0.1(-4.8)=1.08$ & $(1.08-3)^{2}=3.6864$\\
2 & $2(1.08-3)=-3.84$ & $1.08-0.1(-3.84)=1.464$ & $(1.464-3)^{2}=2.3649$\\
\end{tabular}
\end{center}
Since $g\equiv0$, the proximal step reduces to the identity, i.e., $\prox_{\eta g}=I$.

---------------------------------------------------------------------

\section*{Assumptions}
\begin{assumption}[Smoothness of $f$]\label{as:lip}
$f$ is convex and differentiable with $L$‑Lipschitz gradient:
\[
\|\grad f(x)-\grad f(y)\|\le L\|x-y\|,\quad\forall x,y\in\R^{d}.
\]
\end{assumption}
\begin{assumption}[Convexity of $g$]\label{as:conv}
$g$ is proper, closed and convex (no smoothness required).
\end{assumption}
\begin{assumption}[Stepsize]\label{as:step}
The stepsize satisfies $0<\eta\le 1/L$.
\end{assumption}

---------------------------------------------------------------------

\section*{Main Convergence Theorem}
\begin{theorem}\label{thm:rate}
Let $\{x_{k}\}$ be generated by \eqref{eq:update} under Assumptions \ref{as:lip}–\ref{as:step}. Then for any minimizer $x^{\star}\in\arg\min F$,
\[
F(x_{k})-F(x^{\star})
\;\le\;
\frac{\|x_{0}-x^{\star}\|^{2}}{2\eta\,k},
\qquad\forall k\ge 1.
\]
Consequently $F(x_{k})\to F(x^{\star})$ with rate $O(1/k)$.
\end{theorem}

\section*{Theoretical Properties}
\begin{itemize}
  \item \textbf{Stability.} The method is non‑expansive in the sense that the distance to any fixed point does not increase (see Lemma~\ref{lem:nonexp}).
  \item \textbf{Hyperparameter sensitivity.} The bound requires $\eta\le 1/L$; any smaller $\eta$ yields a looser constant $1/(2\eta)$ but preserves convergence.
  \item \textbf{Comparison.} When $g\equiv0$, the method reduces to gradient descent with the classical $O(1/k)$ rate. When $f\equiv0$, the iterates are simply the proximal point algorithm, also enjoying the same rate.
\end{itemize}

\section*{Discussion}
The $O(1/k)$ rate is optimal for first‑order methods on the class of convex composite problems without additional curvature (e.g., strong convexity). Accelerated variants (e.g., FISTA) improve the rate to $O(1/k^{2})$, but the present analysis isolates the role of the proximal operator’s non‑expansiveness.

---------------------------------------------------------------------

\section*{Preliminaries (Definitions and Lemmas)}
\begin{lemma}[Descent Lemma for $f$]\label{lem:descent}
If $\grad f$ is $L$‑Lipschitz, then for any $x,y\in\R^{d}$,
\[
f(y)\le f(x)+\langle\grad f(x),y-x\rangle+\frac{L}{2}\|y-x\|^{2}.
\]
\end{lemma}
\begin{proof}
Standard consequence of Lipschitz continuity of the gradient (integrate the gradient along the line segment). \hfill\qedsymbol
\end{proof}
\begin{lemma}[Non‑expansiveness of the proximal operator]\label{lem:nonexp}
For any $\eta>0$ and any $u,v\in\R^{d}$,
\[
\bigl\|\prox_{\eta g}(u)-\prox_{\eta g}(v)\bigr\|\le\|u-v\|.
\]
\end{lemma}
\begin{proof}
From the optimality condition of the proximal map:
\[
0\in\partial g(\prox_{\eta g}(u))+\frac{1}{\eta}\bigl(\prox_{\eta g}(u)-u\bigr),
\]
and similarly for $v$. Subtract the two inclusions, take inner product with the difference of the two proximal points, and use monotonicity of $\partial g$. This yields the claimed inequality. \hfill\qedsymbol
\end{proof}

---------------------------------------------------------------------

\section*{Main Proof (Step‑by‑Step Derivation)}
We prove Theorem~\ref{thm:rate}. Throughout the proof let $x^{\star}\in\arg\min F$.

\noindent\textbf{Step 1 (Optimality condition of the proximal step).}  
From the definition of $\prox_{\eta g}$,
\[
x_{k+1}= \prox_{\eta g}\!\bigl(x_{k}-\eta\grad f(x_{k})\bigr)
\quad\Longleftrightarrow\quad
0\in \partial g(x_{k+1})+\frac{1}{\eta}\bigl(x_{k+1}-x_{k}+\eta\grad f(x_{k})\bigr).
\tag{1}
\]
Hence there exists $s_{k+1}\in\partial g(x_{k+1})$ such that
\[
s_{k+1}= -\frac{1}{\eta}\bigl(x_{k+1}-x_{k}\bigr)-\grad f(x_{k}).
\tag{2}
\]

\noindent\textbf{Step 2 (Convexity of $g$).}  
By convexity of $g$,
\[
g(x^{\star})\ge g(x_{k+1})+\langle s_{k+1},x^{\star}-x_{k+1}\rangle.
\tag{3}
\]

\noindent\textbf{Step 3 (Descent lemma for $f$).}  
Using Lemma~\ref{lem:descent} with $x=x_{k}$ and $y=x_{k+1}$,
\[
f(x_{k+1})\le f(x_{k})+\langle\grad f(x_{k}),x_{k+1}-x_{k}\rangle+\frac{L}{2}\|x_{k+1}-x_{k}\|^{2}.
\tag{4}
\]

\noindent\textbf{Step 4 (Combine (2)–(4)).}  
Add $f(x_{k+1})$ to both sides of (3) and substitute $s_{k+1}$ from (2):
\[
\begin{aligned}
F(x_{k+1})-F(x^{\star})
&= f(x_{k+1})-f(x^{\star})+g(x_{k+1})-g(x^{\star})\\
&\le f(x_{k})-f(x^{\star})+\langle\grad f(x_{k}),x_{k+1}-x_{k}\rangle
      +\frac{L}{2}\|x_{k+1}-x_{k}\|^{2}\\
&\quad +\langle s_{k+1},x_{k+1}-x^{\star}\rangle .
\end{aligned}
\tag{5}
\]

\noindent\textbf{Step 5 (Insert $s_{k+1}$).}  
Replace $s_{k+1}$ using (2):
\[
\langle s_{k+1},x_{k+1}-x^{\star}\rangle
= -\frac{1}{\eta}\langle x_{k+1}-x_{k},x_{k+1}-x^{\star}\rangle
  -\langle\grad f(x_{k}),x_{k+1}-x^{\star}\rangle .
\tag{6}
\]

\noindent\textbf{Step 6 (Cancel gradient inner products).}  
Plug (6) into (5). The terms $\langle\grad f(x_{k}),x_{k+1}-x_{k}\rangle$ and $-\langle\grad f(x_{k}),x_{k+1}-x^{\star}\rangle$ combine to $-\langle\grad f(x_{k}),x_{k}-x^{\star}\rangle$:
\[
\begin{aligned}
F(x_{k+1})-F(x^{\star})
&\le f(x_{k})-f(x^{\star})-\langle\grad f(x_{k}),x_{k}-x^{\star}\rangle\\
&\quad -\frac{1}{\eta}\langle x_{k+1}-x_{k},x_{k+1}-x^{\star}\rangle
      +\frac{L}{2}\|x_{k+1}-x_{k}\|^{2}.
\end{aligned}
\tag{7}
\]

\noindent\textbf{Step 7 (Convexity of $f$).}  
Convexity of $f$ yields
\[
f(x_{k})-f(x^{\star})\le\langle\grad f(x_{k}),x_{k}-x^{\star}\rangle .
\tag{8}
\]
Insert (8) into (7); the first two terms cancel, leaving
\[
F(x_{k+1})-F(x^{\star})
\le -\frac{1}{\eta}\langle x_{k+1}-x_{k},x_{k+1}-x^{\star}\rangle
    +\frac{L}{2}\|x_{k+1}-x_{k}\|^{2}.
\tag{9}
\]

\noindent\textbf{Step 8 (Algebraic manipulation).}  
Rewrite the inner product:
\[
\langle x_{k+1}-x_{k},x_{k+1}-x^{\star}\rangle
= \tfrac12\Bigl(\|x_{k+1}-x^{\star}\|^{2}
               -\|x_{k}-x^{\star}\|^{2}
               +\|x_{k+1}-x_{k}\|^{2}\Bigr).
\tag{10}
\]
Substituting (10) into (9) gives
\[
\begin{aligned}
F(x_{k+1})-F(x^{\star})
&\le -\frac{1}{2\eta}\Bigl(\|x_{k+1}-x^{\star}\|^{2}
               -\|x_{k}-x^{\star}\|^{2}
               +\|x_{k+1}-x_{k}\|^{2}\Bigr)
   +\frac{L}{2}\|x_{k+1}-x_{k}\|^{2}\\[4pt]
&= \frac{1}{2\eta}\bigl(\|x_{k}-x^{\star}\|^{2}
               -\|x_{k+1}-x^{\star}\|^{2}\bigr)
   -\Bigl(\frac{1}{2\eta}-\frac{L}{2}\Bigr)\|x_{k+1}-x_{k}\|^{2}.
\end{aligned}
\tag{11}
\]

\noindent\textbf{Step 9 (Stepsize condition).}  
By Assumption~\ref{as:step}, $0<\eta\le 1/L$, thus $\frac{1}{\eta}\ge L$ and the coefficient of $\|x_{k+1}-x_{k}\|^{2}$ in (11) is non‑negative:
\[
\frac{1}{2\eta}-\frac{L}{2}\ge 0.
\]
Hence we can drop the non‑negative term to obtain the fundamental inequality
\[
F(x_{k+1})-F(x^{\star})\le
\frac{1}{2\eta}\bigl(\|x_{k}-x^{\star}\|^{2}
               -\|x_{k+1}-x^{\star}\|^{2}\bigr).
\tag{12}
\]
\textbf{[Step 12]}

\noindent\textbf{Step 10 (Telescoping sum).}  
Sum \eqref{12} from $k=0$ to $k=K-1$:
\[
\sum_{k=0}^{K-1}\bigl(F(x_{k+1})-F(x^{\star})\bigr)
\le\frac{1}{2\eta}\bigl(\|x_{0}-x^{\star}\|^{2}-\|x_{K}-x^{\star}\|^{2}\bigr)
\le\frac{\|x_{0}-x^{\star}\|^{2}}{2\eta}.
\tag{13}
\]
Since each term $F(x_{k+1})-F(x^{\star})$ is non‑negative, the smallest of them is bounded by the average:
\[
F(x_{K})-F(x^{\star})
\le\frac{1}{K}\sum_{k=0}^{K-1}\bigl(F(x_{k+1})-F(x^{\star})\bigr)
\le\frac{\|x_{0}-x^{\star}\|^{2}}{2\eta\,K}.
\tag{14}
\]
\textbf{[Step 14]}

\noindent\textbf{Conclusion.} Equation \eqref{14} is precisely the claim of Theorem~\ref{thm:rate}. ∎

---------------------------------------------------------------------

\section*{Rate Derivation (With Constants)}
From \eqref{14} we have an explicit constant:
\[
\boxed{\;
F(x_{k})-F(x^{\star})\le \frac{R^{2}}{2\eta\,k},
\qquad R:=\|x_{0}-x^{\star}\|.
\;}
\]
If the stepsize is chosen as the maximal admissible $\eta=1/L$, the bound becomes
\[
F(x_{k})-F(x^{\star})\le \frac{L\,R^{2}}{2k}.
\]

\section*{Special Cases}
\begin{itemize}
  \item \textbf{Pure gradient descent ($g\equiv0$).} The proximal operator reduces to the identity; the proof collapses to the classical $O(1/k)$ result for smooth convex functions.
  \item \textbf{Proximal point algorithm ($f\equiv0$).} The update is $x_{k+1}=\prox_{\eta g}(x_{k})$, and the same derivation yields $g(x_{k})-g(x^{\star})\le\|x_{0}-x^{\star}\|^{2}/(2\eta k)$.
  \item \textbf{Strongly convex $F$.} If $F$ is $\mu$‑strongly convex, a straightforward modification of Step~9 gives a linear rate $O\bigl((1-\eta\mu)^{k}\bigr)$.
\end{itemize}

\end{document}